\documentclass[conference]{IEEEtran}
\usepackage{cite}
\usepackage{amsmath,amssymb,amsfonts}
\usepackage{graphicx}
\IEEEoverridecommandlockouts
\usepackage{amsmath}
\usepackage{algorithmic}
\usepackage{graphicx}
\usepackage{textcomp}
\usepackage{xcolor}
\usepackage{float}
\usepackage{amsfonts}
\usepackage{amssymb}
\usepackage{graphicx}
\usepackage{cite}
\usepackage{url}
\usepackage{placeins}
\usepackage{comment}
\usepackage{hyperref}
\usepackage{tikz}
\usetikzlibrary{positioning}

\begin{document}

\title{Título de Nuestro Artículo en Formato IEEE}

\author{\IEEEauthorblockN{Autor Uno, Autor Dos, Autor Tres}
\IEEEauthorblockA{\textit{Departamento de Ingeniería} \\
\textit{Universidad XYZ}\\
Ciudad, País}}

\maketitle

\begin{abstract}
Este es el resumen de nuestro artículo configurado con el formato IEEE.
\end{abstract}

\begin{IEEEkeywords}
Palabras clave, LaTeX, Linux, Git.
\end{IEEEkeywords}

\section{Introducción}
Aquí empieza la introducción del documento.

\end{document}
\documentclass[conference]{IEEEtran}
\usepackage{cite}
\usepackage{amsmath,amssymb,amsfonts}
\usepackage{graphicx}

\begin{document}

\title{Título de Nuestro Artículo en Formato IEEE}

\author{\IEEEauthorblockN{George Briceño Celestino}
\IEEEauthorblockA{\textit{Escuela de Ingeniería Electrónica} \\
\textit{Instituto Tecnológico de Costa Rica}\\
Cartago, Costa Rica \\
rousell@estudiantec.cr}
\and

\IEEEauthorblockN{Laura Elena Brenes Espinoza}
\IEEEauthorblockA{\textit{Escuela de Ingeniería Electrónica} \\
\textit{Instituto Tecnológico de Costa Rica}\\
Alajuela, Costa Rica \\
lauritabrenes@estudiantec.cr}
\and

\IEEEauthorblockN{Marlon Méndez Naranjo}
\IEEEauthorblockA{\textit{Escuela de Ingeniería Electrónica} \\
\textit{Instituto Tecnológico de Costa Rica}\\
Cartago, Costa Rica \\
mnaranjo@estudiantec.cr}
\and

\IEEEauthorblockN{Rodrigo Venegas Mora}
\IEEEauthorblockA{\textit{Escuela de Ingeniería Electrónica} \\
\textit{Instituto Tecnológico de Costa Rica}\\
Alajuela, Costa Rica \\
rovenegas@estudiantec.cr}
}

\maketitle

\begin{abstract}
Este es el resumen de nuestro artículo configurado con el formato IEEE.
\end{abstract}

\begin{IEEEkeywords}
Palabras clave, LaTeX, Linux, Git.
\end{IEEEkeywords}

\section{Introducción}
Aquí empieza la introducción del documento.

\end{document}